\documentclass{ieeeaccess}

\usepackage{cite}
\usepackage{sectsty}
\usepackage{float}
\usepackage{url}
\usepackage{multirow}
\usepackage{algorithm}
\usepackage{algpseudocode}
\usepackage{amsmath}

\usepackage[colorlinks=true, urlcolor=blue, linkcolor=blue, citecolor=blue]{hyperref}
\usepackage{graphicx}
\usepackage{textcomp}
\usepackage{xcolor}
\usepackage{booktabs}
\definecolor{orcidgreen}{HTML}{A6CE39}
\usepackage{tabularx}
\usepackage{multirow}

\newcommand{\orcidicon}{\raisebox{-0.1ex}{\textcolor{orcidgreen}{\textsf{\textbf{\footnotesize \textcircled{\tiny iD}}}}}}
\newcommand{\orcidlink}[1]{\href{https://orcid.org/#1}{\orcidicon}}

\sectionfont{\bfseries}
\subsectionfont{\bfseries}
\subsubsectionfont{\bfseries}

\def\BibTeX{{\rm B\kern-.05em{\sc i\kern-.025em b}\kern-.08em
T\kern-.1667em\lower.7ex\hbox{E}\kern-.125emX}}

\begin{document}

\history{Date of publication xxxx 00, 0000, date of current version xxxx 00, 0000.}
\doi{10.1109/ACCESS.2024.XXXXXXX}

\title{Keystroke Dynamics Authentication via Supervised Contrastive Learning and Per-User Adapters}

\author{
\uppercase{Atul~Kumar~Mishra}\authorrefmark{1}\orcidlink{0009-0007-9697-0261},
\uppercase{Harsh~Sagar}\authorrefmark{1}\orcidlink{0009-0006-6097-2479},
\uppercase{Amitnash~Gupta}\authorrefmark{1}\orcidlink{0009-0004-8299-5546},
\uppercase{Dr.~Suma~B}\authorrefmark{2}\orcidlink{0000-0001-7337-1855},
and
\uppercase{Neethu~S}\authorrefmark{3}\orcidlink{0000-0002-8440-5229}
}

\address[1]{Department of Computer Science and Engineering, 
Rashtreeya Vidyalaya College of Engineering (RVCE), 
Bengaluru, Karnataka, India.}

\address[2]{Department of Computer Science and Engineering, 
Rashtreeya Vidyalaya College of Engineering (RVCE), 
Bengaluru, Karnataka, India.}

\address[3]{Department of Computer Science and Engineering, 
Rashtreeya Vidyalaya College of Engineering (RVCE), 
Bengaluru, Karnataka, India.}

\corresp{Corresponding author: Atul~Kumar~Mishra (e-mail: atulkumarmishra.cy23@rvce.edu.in).}

\markboth
{Atul Kumar Mishra \headeretal: Keystroke Dynamics Authentication via Supervised Contrastive Learning and Per-User Adapters}
{Atul Kumar Mishra \headeretal: Keystroke Dynamics Authentication via Supervised Contrastive Learning and Per-User Adapter



\begin{abstract}
We propose a two-stage framework for keystroke dynamics authentication that combines supervised contrastive pretraining with per-user adapter fine-tuning. In the first stage, a global 1D Convolutional Neural Network (CNN) encoder is trained using Supervised Contrastive Loss (SupCon) on the CMU Keystroke Dynamics Benchmark, learning generalizable keystroke embeddings that achieve an Area Under the Curve (AUC) of 0.90 for cross-user verification on unseen users. In the second stage, after user enrollment, a lightweight Multi-Layer Perceptron (MLP) adapter with approximately 16,000 parameters is trained on the user's enrollment samples to personalize the embedding space. This per-user adapter reduces the Equal Error Rate (EER) from 3.3\% (encoder-only with centroid matching) to 2.6\% for enrolled users. Our key insight is that per-user adapters serve a fundamentally different role than the global encoder: adapters are personalization modules effective only after enrollment, not generalization mechanisms designed for cross-user transfer. This separation enables both fast enrollment requiring only 5--10 keystroke samples and strong verification performance for continuous authentication scenarios. To the best of our knowledge, this is the first application of supervised contrastive learning to keystroke dynamics and the first two-stage framework explicitly separating generalization from personalization in behavioral biometrics.
\end{abstract}

\begin{keywords}
Behavioral biometrics, contrastive learning, continuous authentication, keystroke dynamics, per-user personalization, user authentication.
\end{keywords}

\titlepgskip=-15pt

\maketitle

\section{Introduction}
\label{sec:introduction}
\PARstart{K}{eystroke} dynamics is a behavioral biometric modality that analyzes the unique typing patterns of individuals for authentication purposes. Unlike physiological biometrics such as fingerprints or iris scans, keystroke dynamics offers non-intrusive continuous authentication without requiring specialized hardware~\cite{b1,b2}. Users can be verified transparently during normal computer interaction, making it particularly suitable for continuous authentication scenarios where identity must be maintained throughout a session~\cite{b3}.

Traditional approaches to keystroke dynamics authentication rely on statistical models or classical machine learning algorithms such as Random Forests and Support Vector Machines~\cite{b4,b5}. While these methods achieve reasonable performance on benchmark datasets, they suffer from two fundamental limitations. First, they require handcrafted feature engineering that may not capture the full complexity of typing behavior. Second, they do not effectively leverage the shared structure across users that could improve generalization~\cite{b6}.

Recent advances in deep learning have shown promising results for keystroke dynamics, with Convolutional Neural Networks (CNNs) and Recurrent Neural Networks (RNNs) learning representations directly from raw timing data~\cite{b7,b8}. However, most existing deep learning approaches train models for user identification (closed-set classification) rather than verification (open-set authentication), limiting their practical applicability~\cite{b9}.

In this paper, we propose a novel two-stage framework that addresses these limitations:

\begin{enumerate}
\item \textbf{Supervised Contrastive Pretraining:} We train a 1D-CNN encoder using Supervised Contrastive Loss (SupCon)~\cite{b10} to learn embeddings where samples from the same user are clustered together while samples from different users are pushed apart. This produces generalizable representations that work even for users not seen during training.

\item \textbf{Per-User Adapter Fine-tuning:} After user enrollment, we train a lightweight MLP adapter ($\sim$16K parameters) on the user's enrollment samples. This personalizes the embedding space for the specific user, reducing EER from 3.3\% to 2.6\%.
\end{enumerate}

Our key insight is that these two components serve fundamentally different roles: the encoder provides \textit{generalization} across users, while adapters provide \textit{personalization} for enrolled users. This separation enables fast enrollment (5--10 samples) with strong verification performance.

The main contributions of this work are:
\begin{itemize}
\item First application of Supervised Contrastive Learning to keystroke dynamics authentication.
\item A two-stage framework explicitly separating generalization (encoder) from personalization (adapter).
\item Comprehensive evaluation on the CMU Keystroke Dynamics Benchmark demonstrating state-of-the-art performance.
\item Analysis of why per-user adapters are effective for enrolled users but do not generalize across users.
\end{itemize}

\section{Related Work}
\label{sec:related}

\subsection{Keystroke Dynamics Authentication}
Keystroke dynamics has been studied extensively as a behavioral biometric since the pioneering work of Monrose and Rubin~\cite{b1}. Early approaches used statistical distance measures such as Euclidean distance and Mahalanobis distance to compare typing patterns against stored templates~\cite{b2}. Killourhy and Maxion~\cite{b4} established the CMU Keystroke Dynamics Benchmark and systematically compared anomaly detection algorithms, finding that the Manhattan (scaled) detector achieved the best performance with an EER of 9.6\%.

Machine learning approaches have gradually replaced statistical methods. Bergadano et al.~\cite{b11} used decision trees and achieved 96\% accuracy on a custom dataset. More recent work has applied ensemble methods such as Random Forests and Gradient Boosting, achieving improved accuracy on benchmark datasets~\cite{b5,b12}.

\subsection{Deep Learning for Keystroke Dynamics}
Deep learning approaches have shown promising results for keystroke dynamics. Deng and Zhong~\cite{b7} applied deep neural networks to keystroke data, demonstrating improved performance over traditional methods. Acien et al.~\cite{b13} proposed TypeNet, a deep learning architecture specifically designed for keystroke biometrics that uses LSTM layers to model temporal dependencies.

Lu et al.~\cite{b8} combined CNNs and RNNs for free-text keystroke authentication, achieving strong results on continuous authentication tasks. Roy et al.~\cite{b14} explored various deep architectures and found that 1D-CNNs are particularly effective for fixed-text keystroke data due to their ability to capture local temporal patterns.

\subsection{Contrastive Learning}
Contrastive learning has emerged as a powerful paradigm for learning representations without explicit labels. Chen et al.~\cite{b15} introduced SimCLR, which uses a contrastive loss to learn visual representations from unlabeled images. Khosla et al.~\cite{b10} extended this to the supervised setting with Supervised Contrastive Loss (SupCon), which leverages label information to create tighter clusters per class.

In the biometrics domain, contrastive learning has been applied to face recognition with great success. FaceNet~\cite{b16} uses triplet loss to learn face embeddings, while more recent work has adopted contrastive objectives for improved performance~\cite{b17}. However, to the best of our knowledge, supervised contrastive learning has not been applied to keystroke dynamics.

\subsection{Deep Metric Learning}
Deep metric learning aims to learn embedding spaces where semantic similarity is reflected by geometric proximity. Beyond triplet loss~\cite{b16}, several advanced loss functions have been proposed. ProxyNCA~\cite{b23} uses learnable proxies as substitutes for class representatives, enabling efficient batch training. ArcFace~\cite{b24} introduces an additive angular margin penalty to enhance intra-class compactness and inter-class discrepancy, achieving state-of-the-art results in face recognition.

Center Loss~\cite{b25} jointly minimizes cross-entropy and the distance between deep features and their corresponding class centers. Multi-Similarity Loss~\cite{b26} considers multiple similarity types within pairs to mine informative samples more effectively. These advances in deep metric learning have primarily focused on visual recognition tasks, with limited exploration in behavioral biometrics.

\subsection{Few-Shot Learning and Adapters}
Few-shot learning addresses the challenge of learning from limited examples. Prototypical Networks~\cite{b18} learn a metric space where classification is performed by computing distances to class prototypes. Siamese Networks~\cite{b19} learn to compare pairs of samples directly. Matching Networks~\cite{b27} combine attention mechanisms with episode-based training for improved few-shot generalization.

Adapter modules have been proposed in the natural language processing domain as lightweight alternatives to full fine-tuning~\cite{b20}. These small bottleneck layers (typically 0.5--5\% of model parameters) are inserted into pretrained models and trained on downstream tasks while keeping the base model frozen. Recent work has extended adapters to vision models~\cite{b28} and introduced low-rank adaptation (LoRA)~\cite{b29} which decomposes weight updates into low-rank matrices.

We adapt this concept to keystroke dynamics, using per-user adapters to personalize a frozen encoder. Unlike NLP adapters that specialize to tasks, our adapters specialize to individual users, creating personalized authentication boundaries in the embedding space.

\subsection{Biometric Template Protection}
Securing stored biometric templates is critical for deployment. Cancelable biometrics~\cite{b30} applies irreversible transformations to templates, allowing revocation if compromised. Fuzzy commitment schemes~\cite{b31} combine biometric data with cryptographic keys. While our current work does not implement template protection, the adapter architecture naturally supports this---adapter weights can be encrypted or stored on secure hardware independently of the shared encoder.

\section{Dataset and Feature Extraction}
\label{sec:dataset}

\subsection{CMU Keystroke Dynamics Benchmark}
We use the CMU Keystroke Dynamics Benchmark~\cite{b4}, the most widely-used public dataset for evaluating fixed-text keystroke authentication systems. The dataset was collected at Carnegie Mellon University under controlled laboratory conditions to ensure data quality.

\subsubsection{Collection Protocol}
The dataset contains 20,400 keystroke samples from 51 participants (33 male, 18 female, ages 18--70). Each participant typed the password ``.tie5Roanl'' 400 times across 8 data collection sessions, with approximately 50 repetitions per session. Sessions were separated by at least one day to capture natural variation in typing behavior over time. Participants used the same keyboard throughout all sessions to control for hardware differences.

\subsubsection{Password Design}
The 11-character password ``.tie5Roanl'' was specifically designed to include:
\begin{itemize}
\item Special characters (period and digit)
\item Mixed case letters requiring shift key usage
\item Common letter combinations and uncommon sequences
\item Sufficient length to capture stable timing patterns
\end{itemize}

This design ensures that typing patterns are distinctive and capture a range of motor control behaviors including finger transitions, shift key coordination, and rhythm.

\subsubsection{Timing Measurements}
For each keystroke event, the dataset records high-resolution timestamps (millisecond precision) for both key press (down) and key release (up) events. From these raw timestamps, three types of derived timing features are computed for each key position.

\subsection{Feature Representation}
We extract three types of timing features for each key in the password:

\begin{itemize}
\item \textbf{Hold Time (H):} Duration between key press and key release for a single key.
\item \textbf{Down-Down Time (DD):} Time between consecutive key presses.
\item \textbf{Up-Down Time (UD):} Time between releasing one key and pressing the next.
\end{itemize}

For the 11-character password, we obtain 11 hold times, 10 down-down times, and 10 up-down times, totaling 31 raw timing features. We organize these into a canonical sequence format of shape $(N, 11, 3)$, where $N$ is the batch size, 11 is the number of key tokens, and 3 represents the features per key (H, DD, UD).

\subsection{Data Preprocessing}
We apply z-score normalization to all timing features using statistics computed from the training set. This ensures that all features are on a comparable scale and helps with neural network optimization. We also apply data augmentation during contrastive training, including:

\begin{itemize}
\item \textbf{Gaussian Noise:} Adding small random perturbations to timing values.
\item \textbf{Time Warping:} Slight stretching or compression of timing intervals.
\end{itemize}

These augmentations create positive pairs for contrastive learning while preserving the essential typing characteristics.

\begin{figure}[t]
  \centering
  \includegraphics[width=\columnwidth]{data_distribution.png}
  \caption{Distribution of typing samples across users in the CMU Keystroke Dynamics Benchmark. Each user contributed 400 samples across 8 sessions.}
  \label{fig:data_dist}
\end{figure}

\section{Proposed Methodology}
\label{sec:methodology}

Our framework consists of two stages: (1) supervised contrastive pretraining of a global encoder, and (2) per-user adapter fine-tuning after enrollment. Fig.~\ref{fig:architecture} illustrates the overall architecture.

\begin{figure}[t]
  \centering
  \includegraphics[width=\columnwidth]{un2.png}
  \caption{Overview of the proposed two-stage framework. \textbf{Stage 1}: A 1D-CNN encoder is pretrained with Supervised Contrastive Loss on all training users. \textbf{Stage 2}: Per-user adapters are trained on enrollment samples while the encoder remains frozen.}
  \label{fig:architecture}
\end{figure}


\subsection{Supervised Contrastive Encoder}
The encoder is a 1D Convolutional Neural Network that processes keystroke sequences of shape $(11, 3)$ and produces 128-dimensional L2-normalized embeddings.

\subsubsection{Architecture}
The encoder consists of three convolutional blocks followed by global average pooling and a fully-connected layer:

\begin{itemize}
\item \textbf{Conv Block 1:} Conv1D($3 \rightarrow 64$, kernel=3) $\rightarrow$ BatchNorm $\rightarrow$ ReLU $\rightarrow$ MaxPool
\item \textbf{Conv Block 2:} Conv1D($64 \rightarrow 128$, kernel=3) $\rightarrow$ BatchNorm $\rightarrow$ ReLU $\rightarrow$ MaxPool
\item \textbf{Conv Block 3:} Conv1D($128 \rightarrow 256$, kernel=3) $\rightarrow$ BatchNorm $\rightarrow$ ReLU $\rightarrow$ GlobalAvgPool
\item \textbf{Embedding:} Linear($256 \rightarrow 128$) $\rightarrow$ L2-Normalize
\end{itemize}

The total number of parameters is approximately 158,000.

\subsubsection{Supervised Contrastive Loss}
We train the encoder using Supervised Contrastive Loss (SupCon)~\cite{b10}, which extends the self-supervised NT-Xent loss to leverage label information:

\begin{equation}
\mathcal{L}_{SupCon} = \sum_{i} \frac{-1}{|P(i)|} \sum_{p \in P(i)} \log \frac{\exp(z_i \cdot z_p / \tau)}{\sum_{a \neq i} \exp(z_i \cdot z_a / \tau)}
\label{eq:supcon}
\end{equation}

where $P(i)$ is the set of indices of samples with the same label as sample $i$, $z_i$ is the L2-normalized embedding of sample $i$, and $\tau = 0.07$ is the temperature parameter.

Unlike cross-entropy loss, SupCon groups \textit{all} samples from the same user as positives, creating tighter per-user clusters in the embedding space.

\subsubsection{Training Details}
We train the encoder for 100 epochs using AdamW optimizer with learning rate $10^{-3}$, weight decay $10^{-4}$, and cosine annealing learning rate schedule. We use a batch size of 128 and apply mixed-precision training for efficiency.

\subsection{Per-User Adapter Design}
After encoder pretraining, we freeze the encoder weights and train a lightweight adapter for each enrolled user.

\subsubsection{Architecture}
The adapter is a simple MLP with residual connection:

\begin{equation}
\text{Adapter}(x) = \text{Linear}_2(\text{ReLU}(\text{Linear}_1(x)))
\label{eq:adapter}
\end{equation}

where $\text{Linear}_1: 128 \rightarrow 64$ and $\text{Linear}_2: 64 \rightarrow 128$.

We found that \textbf{removing the residual connection} (i.e., not adding the input back to the output) improves performance, as it allows the adapter to fully transform the embedding rather than simply refining it.

The total number of parameters per adapter is approximately 16,000, enabling efficient storage and fast training.

\subsubsection{Training Objective}
We train adapters using triplet loss with hard negative mining:

\begin{equation}
\mathcal{L}_{triplet} = \max(0, d(a, p) - d(a, n) + m)
\label{eq:triplet}
\end{equation}

where $a$ is the anchor (user centroid), $p$ is a positive sample (same user), $n$ is a hard negative (most similar impostor), and $m = 0.3$ is the margin.

\subsubsection{Enrollment Protocol}
Given $k$ enrollment samples from a user:
\begin{enumerate}
\item Compute encoder embeddings for all samples.
\item Initialize adapter with random weights.
\item Train adapter for 100 steps using triplet loss.
\item Compute centroid from adapted embeddings.
\item Store adapter weights and centroid as user template.
\end{enumerate}

\subsection{Verification Pipeline}
For verification, given a query keystroke sample and claimed identity:

\begin{enumerate}
\item Compute encoder embedding: $z = f_{enc}(x)$
\item Apply user's adapter: $z' = f_{adapt}(z)$
\item Compute similarity: $s = \cos(z', c_{user})$
\item Accept if $s \geq \theta_{user}$
\end{enumerate}

The threshold $\theta_{user}$ is calibrated per-user using the EER criterion on enrollment samples.

\section{Experimental Setup}
\label{sec:experiments}

\subsection{Implementation Details}
Our framework is implemented in PyTorch 2.0 with CUDA 11.8 support. All experiments are conducted on a workstation equipped with an NVIDIA RTX 3080 GPU (10 GB VRAM), Intel Core i7-12700K CPU, and 32 GB RAM. Training uses automatic mixed precision (AMP) via PyTorch's GradScaler to reduce memory consumption and accelerate computation.

Table~\ref{tab:hyperparams} summarizes the complete hyperparameter configuration for both encoder pretraining and adapter fine-tuning stages.

\begin{table}[h]
\caption{Complete Hyperparameter Configuration}
\label{tab:hyperparams}
\centering
\begin{tabular}{|l|l|c|}
\hline
\textbf{Stage} & \textbf{Hyperparameter} & \textbf{Value} \\
\hline
\multirow{10}{*}{Encoder} & Embedding dimension & 128 \\
& Conv channels & [64, 128, 256] \\
& Kernel size & 3 \\
& Batch size & 128 \\
& Epochs & 100 \\
& Learning rate & $10^{-3}$ \\
& Weight decay & $10^{-4}$ \\
& Optimizer & AdamW \\
& LR scheduler & Cosine annealing \\
& Temperature $\tau$ & 0.07 \\
\hline
\multirow{8}{*}{Adapter} & Hidden dimension & 64 \\
& Dropout rate & 0.2 \\
& Learning rate & $10^{-3}$ \\
& Optimizer & Adam \\
& Max training steps & 200 \\
& Triplet margin & 0.3 \\
& Hard negative mining & Yes \\
& Residual connection & No \\
\hline
\multirow{4}{*}{Augmentation} & Jitter probability & 0.5 \\
& Jitter scale & $\pm$3\% \\
& Time warp probability & 0.2 \\
& Time warp scale & $\pm$10\% \\
\hline
\end{tabular}
\end{table}

\subsection{Evaluation Protocols}
We evaluate our system using two complementary protocols designed to measure different aspects of system performance.

\subsubsection{Cross-User Evaluation (User-Disjoint Split)}
We partition users into disjoint training (36 users, 70\%), validation (7 users, 14\%), and test (8 users, 16\%) sets. The encoder is trained exclusively on training users and evaluated on completely unseen test users. This protocol measures the \textit{generalization} capability of learned representations to users never encountered during training---a critical requirement for practical deployment where new users must be authenticated without retraining.

\subsubsection{In-Distribution Evaluation (Temporal Split)}
We employ a temporal split where all 51 users are present in both training and test sets, but test samples come from chronologically later sessions. Specifically, sessions 1--6 are used for training/enrollment, while sessions 7--8 are reserved for testing. Adapters are trained on $k$ enrollment samples (where $k \in \{1, 3, 5, 10\}$) and evaluated on held-out samples from the \textit{same} users. This protocol measures \textit{personalization} performance and session-to-session consistency, simulating the realistic scenario of continuous authentication after enrollment.

\subsection{Baselines}
We compare against the following baseline methods:

\begin{itemize}
\item \textbf{Random Forest (RF):} Ensemble of 100 decision trees trained on 47-dimensional aggregated features (31 raw timing features plus 16 statistical aggregates including mean, standard deviation, min, max for each feature group). Uses Gini impurity for splitting and bootstrap aggregation.

\item \textbf{XGBoost:} Gradient boosted trees with 100 estimators, max depth 6, learning rate 0.1, and the same 47-dimensional feature set. Trained with log loss objective for binary verification.

\item \textbf{Encoder-only (Centroid):} Our pretrained 1D-CNN encoder with centroid-based matching. For each user, we compute the mean embedding (centroid) of enrollment samples and use cosine similarity to the query for verification.
\end{itemize}

\subsection{Evaluation Metrics}
We report the following standard biometric verification metrics:

\begin{itemize}
\item \textbf{Equal Error Rate (EER):} The operating point where False Acceptance Rate (FAR) equals False Rejection Rate (FRR). Computed by interpolating the FAR/FRR curves. Lower EER indicates better discrimination.

\item \textbf{Area Under ROC Curve (AUC):} The probability that a randomly chosen genuine pair scores higher than a randomly chosen impostor pair. Ranges from 0.5 (random) to 1.0 (perfect). Threshold-independent measure of overall discriminative performance.

\item \textbf{False Acceptance Rate at 1\% FRR (FAR@1\%FRR):} The impostor acceptance rate when the system rejects only 1\% of legitimate users. Important for high-security applications.
\end{itemize}

\subsection{Statistical Significance}
All experiments involving randomness (enrollment sample selection, train/val splits) are repeated 5 times with different random seeds. We report mean $\pm$ standard deviation across runs. Statistical significance is assessed using paired t-tests with $\alpha = 0.05$.

\section{Results and Analysis}
\label{sec:results}

\subsection{Cross-User Evaluation}

Table~\ref{tab:crossuser} presents results on the user-disjoint split, measuring generalization to unseen users.

\begin{table}[h]
\caption{Cross-User Evaluation Results (User-Disjoint Split)}
\label{tab:crossuser}
\centering
\begin{tabular}{|l|c|c|}
\hline
\textbf{Model} & \textbf{AUC} & \textbf{EER} \\
\hline
Random Forest & 0.80 & 26.3\% \\
XGBoost & 0.66 & 39.4\% \\
\textbf{Encoder (centroid)} & \textbf{0.90} & \textbf{17.7\%} \\
Encoder + Adapter & 0.45 & 50.2\% \\
\hline
\end{tabular}
\end{table}

The encoder-only approach achieves the best cross-user performance with AUC = 0.90 and EER = 17.7\%, substantially outperforming classical baselines. Notably, adapters \textit{degrade} cross-user performance, which is expected since adapters are trained on specific users and cannot generalize to unseen users.

\begin{figure}[t]
  \centering
  \includegraphics[width=\columnwidth]{roc_consolidated.png}
  \caption{ROC curves for cross-user evaluation. The encoder-only approach (blue) substantially outperforms baseline methods.}
  \label{fig:roc_cross}
\end{figure}


\subsection{In-Distribution Evaluation}

Table~\ref{tab:indist} presents results on the temporal split, measuring personalization performance for enrolled users.

\begin{table}[h]
\caption{In-Distribution Evaluation Results (Temporal Split)}
\label{tab:indist}
\centering
\begin{tabular}{|l|c|c|}
\hline
\textbf{Model} & \textbf{AUC} & \textbf{EER} \\
\hline
Encoder (centroid) & 0.994 & 3.3\% \\
Encoder + Adapter (k=5) & 0.993 & 2.8\% \\
\textbf{Encoder + Adapter (k=10)} & \textbf{0.991} & \textbf{2.6\%} \\
\hline
\end{tabular}
\end{table}

When evaluated on enrolled users, adapters provide significant improvement, reducing EER from 3.3\% to 2.6\% (21\% relative improvement). This demonstrates that adapters are effective for personalization after enrollment.

\begin{figure}[t]
  \centering
  \includegraphics[width=\columnwidth]{eer_vs_shots.png}
  \caption{EER as a function of enrollment shots. Adapter-based methods consistently outperform centroid-only matching.}
  \label{fig:eer_shots}
\end{figure}


\subsection{Ablation Studies}

We conduct comprehensive ablation studies to understand the contribution of each component and the sensitivity to hyperparameter choices.

\subsubsection{Adapter Architecture Ablation}
We compare adapter variants with and without residual connection:

\begin{table}[h]
\caption{Adapter Architecture Ablation (k=10 shots)}
\label{tab:ablation}
\centering
\begin{tabular}{|l|c|c|c|}
\hline
\textbf{Adapter Variant} & \textbf{AUC} & \textbf{EER} & \textbf{Params} \\
\hline
With Residual & 0.993 & 3.4\% & 16,576 \\
\textbf{Without Residual} & \textbf{0.991} & \textbf{2.6\%} & 16,576 \\
No Adapter (centroid only) & 0.994 & 3.3\% & 0 \\
\hline
\end{tabular}
\end{table}

Removing the residual connection improves EER by 0.8 percentage points. We hypothesize that the residual connection preserves too much of the original embedding, limiting the adapter's ability to create user-specific decision boundaries. Without the residual path, the adapter can fully transform the embedding to optimize for the individual user.

\subsubsection{Temperature Sensitivity}
The temperature parameter $\tau$ in the SupCon loss controls the concentration of the embedding distribution. Table~\ref{tab:temp_ablation} shows performance across different temperature values.

\begin{table}[h]
\caption{Temperature Sensitivity Analysis}
\label{tab:temp_ablation}
\centering
\begin{tabular}{|c|c|c|c|}
\hline
\textbf{Temperature $\tau$} & \textbf{Cross-User AUC} & \textbf{In-Dist AUC} & \textbf{In-Dist EER} \\
\hline
0.01 & 0.82 & 0.976 & 5.8\% \\
0.05 & 0.88 & 0.991 & 3.1\% \\
\textbf{0.07} & \textbf{0.90} & \textbf{0.994} & \textbf{3.3\%} \\
0.10 & 0.89 & 0.992 & 3.4\% \\
0.20 & 0.85 & 0.987 & 4.2\% \\
\hline
\end{tabular}
\end{table}

The optimal temperature of $\tau = 0.07$ aligns with values reported in the original SimCLR and SupCon papers. Lower temperatures create harder contrastive objectives but may lead to optimization instability; higher temperatures produce smoother gradients but weaker clustering.

\subsubsection{Embedding Dimension}
We evaluate the impact of embedding dimensionality on both performance and model size:

\begin{table}[h]
\caption{Embedding Dimension Analysis}
\label{tab:embed_dim}
\centering
\begin{tabular}{|c|c|c|c|c|}
\hline
\textbf{Dim} & \textbf{AUC} & \textbf{EER} & \textbf{Encoder Params} & \textbf{Adapter Params} \\
\hline
64 & 0.988 & 3.8\% & 142K & 8,448 \\
\textbf{128} & \textbf{0.994} & \textbf{3.3\%} & 158K & 16,576 \\
256 & 0.993 & 3.4\% & 191K & 33,024 \\
512 & 0.992 & 3.5\% & 257K & 66,048 \\
\hline
\end{tabular}
\end{table}

An embedding dimension of 128 provides the best trade-off between performance and efficiency. Larger dimensions do not improve performance and increase storage requirements. The relatively modest dimensionality is sufficient for 51 users; larger deployments may benefit from higher dimensions.

\subsubsection{Data Augmentation Ablation}
We ablate each augmentation type to understand their individual contributions:

\begin{table}[h]
\caption{Data Augmentation Ablation}
\label{tab:aug_ablation}
\centering
\begin{tabular}{|l|c|c|}
\hline
\textbf{Augmentation Setting} & \textbf{AUC} & \textbf{EER} \\
\hline
No augmentation & 0.981 & 4.5\% \\
Gaussian jitter only & 0.989 & 3.7\% \\
Time warp only & 0.984 & 4.1\% \\
Jitter + Time warp (full) & \textbf{0.994} & \textbf{3.3\%} \\
\hline
\end{tabular}
\end{table}

Gaussian jitter provides the largest improvement, simulating natural variations in key press timing. Time warping contributes additional gains by modeling changes in overall typing speed. The combination of both augmentations achieves the best performance.

\subsection{Enrollment Sample Analysis}

\subsubsection{Effect of Enrollment Shots}
Table~\ref{tab:shots} shows the detailed performance as a function of the number of enrollment samples $k$.

\begin{table}[h]
\caption{Performance vs. Number of Enrollment Shots}
\label{tab:shots}
\centering
\begin{tabular}{|c|c|c|c|c|}
\hline
\textbf{k} & \textbf{Centroid EER} & \textbf{Adapter EER} & \textbf{Improvement} & \textbf{Training Time} \\
\hline
1 & 5.2\% & 5.1\% & 1.9\% & 0.8s \\
3 & 3.9\% & 3.5\% & 10.3\% & 1.2s \\
5 & 3.5\% & 2.8\% & 20.0\% & 1.5s \\
10 & 3.3\% & 2.6\% & 21.2\% & 2.1s \\
\hline
\end{tabular}
\end{table}

With only $k=1$ sample, adapters provide minimal benefit as there is insufficient data to learn user-specific transformations. The adapter advantage grows with more samples, plateauing around $k=10$. We recommend $k=5$ as a practical trade-off between enrollment effort and verification accuracy.

\subsection{Per-User Analysis}

\subsubsection{EER Distribution Across Users}
Not all users are equally easy to authenticate. Fig.~\ref{fig:eer_dist} shows the distribution of per-user EER values across all 51 subjects.

The per-user EER ranges from 0.8\% (easiest user) to 6.2\% (hardest user), with a median of 2.4\% and interquartile range of 1.7\%--3.5\%. This variation reflects inherent differences in typing consistency across individuals.

\subsubsection{Factors Affecting User Difficulty}
We analyze correlations between user characteristics and authentication difficulty:

\begin{table}[h]
\caption{Correlation with Per-User EER}
\label{tab:user_factors}
\centering
\begin{tabular}{|l|c|c|}
\hline
\textbf{User Characteristic} & \textbf{Pearson r} & \textbf{p-value} \\
\hline
Intra-user timing variance & 0.72 & $<$0.001 \\
Mean typing speed & -0.31 & 0.027 \\
Session-to-session consistency & -0.65 & $<$0.001 \\
\hline
\end{tabular}
\end{table}

Users with high intra-user timing variance (inconsistent typists) are harder to authenticate because their enrollment samples do not accurately represent their test behavior. Fast typists tend to be more consistent, leading to lower EER. Session-to-session consistency is the strongest predictor of authentication success.

\subsubsection{Best and Worst Performing Users}
Table~\ref{tab:user_extremes} shows the 5 easiest and 5 hardest users to authenticate.

\begin{table}[h]
\caption{Per-User EER Extremes}
\label{tab:user_extremes}
\centering
\begin{tabular}{|c|c|c||c|c|c|}
\hline
\multicolumn{3}{|c||}{\textbf{Easiest Users}} & \multicolumn{3}{c|}{\textbf{Hardest Users}} \\
\hline
\textbf{User} & \textbf{EER} & \textbf{Variance} & \textbf{User} & \textbf{EER} & \textbf{Variance} \\
\hline
s022 & 0.8\% & 0.012 & s041 & 6.2\% & 0.089 \\
s015 & 1.1\% & 0.015 & s038 & 5.8\% & 0.076 \\
s007 & 1.2\% & 0.018 & s029 & 5.5\% & 0.071 \\
s033 & 1.3\% & 0.016 & s044 & 5.2\% & 0.068 \\
s019 & 1.4\% & 0.019 & s012 & 4.9\% & 0.064 \\
\hline
\end{tabular}
\end{table}

\subsection{Training Dynamics}

\subsubsection{Learning Curve Analysis}
The encoder was trained for 100 epochs with cosine annealing learning rate schedule. Training loss decreased from 3.99 to 1.91 over the training period, as shown in Fig.~\ref{fig:training_curves}. Validation loss plateaued around epoch 8 at approximately 4.97, indicating that the model's generalization capability stabilizes early while training loss continues to decrease.
\begin{figure}[t]
  \centering
  \includegraphics[width=\columnwidth]{training_curves.png}
  \caption{Training and validation loss curves over 100 epochs. The training loss decreases smoothly while the validation loss plateaus early, indicating good generalization.}
  \label{fig:training_curves}
\end{figure}


The gap between training and validation loss suggests mild overfitting to the training users, which is expected in contrastive learning where the model optimizes for user separation in the training set. However, the cross-user evaluation results confirm that the learned representations still generalize well to unseen users.

\subsubsection{Score Distribution Analysis}
For the in-distribution evaluation with $k=10$ enrollment samples, we observe clear separation between genuine and impostor score distributions, as illustrated in Fig.~\ref{fig:score_dist}:

\begin{itemize}
\item \textbf{Genuine pairs:} Mean cosine similarity = 0.845 ($\sigma$ = 0.089)
\item \textbf{Impostor pairs:} Mean cosine similarity = -0.459 ($\sigma$ = 0.182)
\item \textbf{Separation:} 14.6 standard deviations (d' = 14.6)
\end{itemize}

\begin{figure}[t]
  \centering
  \includegraphics[width=\columnwidth]{score_distributions.png}
  \caption{Distribution of cosine similarity scores for genuine (same-user) and impostor (different-user) pairs. The clear separation enables low-EER authentication.}
  \label{fig:score_dist}
\end{figure}


The large separation (over 14 standard deviations) explains the low EER and high AUC. The impostor distribution has higher variance, reflecting the diversity of typing patterns across different users.


\subsection{Embedding Visualization}
\begin{figure}[t]
  \centering
  \includegraphics[width=\columnwidth]{tsne_embeddings.png}
  \caption{t-SNE visualization of encoder embeddings. Each color represents a different user. Users form distinct clusters, demonstrating the encoder's discriminative power.}
  \label{fig:tsne}
\end{figure}





Fig.~\ref{fig:tsne} shows t-SNE visualization of encoder embeddings for 10 randomly selected users. Users form clearly separated clusters, indicating that the supervised contrastive objective successfully learns discriminative representations.

\section{Discussion}
\label{sec:discussion}

\subsection{Why Adapters Fail Cross-User but Succeed In-Distribution}

Our results reveal a fundamental distinction between generalization and personalization in biometric systems:

\begin{itemize}
\item \textbf{Cross-user evaluation} tests whether a model can verify users it has never seen before. Adapters fail here because they are \textit{per-user} modules trained on specific enrollment samples. An adapter trained for User A has no mechanism to transfer to User B.

\item \textbf{In-distribution evaluation} tests whether a model can verify \textit{enrolled} users over time. This is the relevant setting for continuous authentication, where each user has their own adapter trained on their enrollment data.
\end{itemize}

This behavior is \textbf{by design, not a failure}. Adapters are personalization modules, not generalization mechanisms.

\subsection{Implications for System Design}

Our framework naturally supports two operational modes:

\begin{enumerate}
\item \textbf{Generalization Mode (Pre-Enrollment):} For new users before enrollment, use the encoder with centroid matching. This provides reasonable authentication (EER = 17.7\%) with minimal setup.

\item \textbf{Personalization Mode (Post-Enrollment):} After collecting 5--10 enrollment samples, train an adapter for the user. This provides much stronger authentication (EER = 2.6\%).
\end{enumerate}

The transition between modes is seamless, requiring only a few additional keystroke samples.

\subsection{Computational Efficiency}

A key advantage of our two-stage approach is computational efficiency at both training and inference time.

\subsubsection{Training Costs}
Table~\ref{tab:training_cost} summarizes the computational requirements for each component.

\begin{table}[h]
\caption{Computational Requirements}
\label{tab:training_cost}
\centering
\begin{tabular}{|l|c|c|c|}
\hline
\textbf{Component} & \textbf{Time} & \textbf{Parameters} & \textbf{Storage} \\
\hline
Encoder training (100 epochs) & 195.6s & 158K & 632 KB \\
Adapter training (per user) & 1.5--2.1s & 16.5K & 66 KB \\
Centroid computation & 0.02s & -- & 512 B \\
\hline
\multicolumn{4}{|l|}{\textit{Total for 51 users:}} \\
\hline
All adapters & $\sim$100s & 842K & 3.4 MB \\
All centroids & $\sim$1s & -- & 26 KB \\
\hline
\end{tabular}
\end{table}

The encoder is trained once and shared across all users, amortizing its cost. Each user requires only 1.5--2 seconds to enroll, enabling real-time user registration.

\subsubsection{Inference Latency}
Verification latency is critical for user experience. We measure end-to-end inference time on both GPU and CPU:

\begin{table}[h]
\caption{Inference Latency per Verification}
\label{tab:latency}
\centering
\begin{tabular}{|l|c|c|}
\hline
\textbf{Component} & \textbf{GPU (RTX 3080)} & \textbf{CPU (i7-12700K)} \\
\hline
Feature extraction & 0.1 ms & 0.3 ms \\
Encoder forward pass & 0.8 ms & 3.2 ms \\
Adapter forward pass & 0.1 ms & 0.4 ms \\
Cosine similarity & 0.01 ms & 0.02 ms \\
\hline
\textbf{Total} & \textbf{1.01 ms} & \textbf{3.92 ms} \\
\hline
\end{tabular}
\end{table}

Even on CPU, verification completes in under 4 milliseconds, well below the threshold for real-time interaction. The lightweight adapter adds only 10--12\% overhead compared to encoder-only verification.

\subsubsection{Scalability Analysis}
We analyze how storage and computation scale with the number of users:

\begin{equation}
\text{Storage} = S_{enc} + n \times (S_{adapter} + S_{centroid})
\label{eq:storage}
\end{equation}

where $S_{enc} = 632$ KB (shared), $S_{adapter} = 66$ KB, and $S_{centroid} = 512$ B.

For 10,000 users: $632 + 10000 \times 66.5 \approx 665$ MB total storage.
For 1,000,000 users: $\approx 65$ GB total storage.

Verification time is independent of the number of users (constant lookup), making the system suitable for large-scale deployment. The per-user storage of 67 KB is comparable to a small image file.

\subsubsection{Memory Efficiency}
During inference, only the encoder (632 KB) and the target user's adapter (66 KB) need to be loaded in memory. Other users' adapters can remain on disk. This enables deployment on memory-constrained edge devices such as smartphones or embedded systems.

\subsection{Security Considerations}

\subsubsection{Replay Attack Resistance}
Keystroke dynamics inherently includes temporal information that varies with each input. Unlike static biometrics, captured keystroke sequences cannot be trivially replayed since the exact timing pattern would need to be reproduced. However, sophisticated attackers might attempt to mimic timing patterns, motivating the need for:
\begin{itemize}
\item Liveness detection via typing rhythm analysis
\item Challenge-response protocols with different passphrases
\item Fusion with other behavioral signals (mouse dynamics, session context)
\end{itemize}

\subsubsection{Template Storage Security}
Our adapter-based architecture provides a natural security boundary:
\begin{itemize}
\item The shared encoder can be stored publicly or on a server
\item Per-user adapters and centroids can be encrypted or stored on secure hardware (TPM, secure enclave)
\item Compromise of one user's adapter does not expose other users
\end{itemize}

If an adapter is compromised, it can be revoked and re-enrolled with new keystroke samples, providing a form of cancelable biometrics.

\subsection{Comparison with Prior Work}

Table~\ref{tab:prior_work} compares our approach with previous keystroke authentication methods on the CMU Benchmark.

\begin{table*}[t]
\caption{Comparison with Prior Work on CMU Keystroke Dynamics Benchmark}
\label{tab:prior_work}
\centering
\begin{tabular}{|l|c|c|c|c|c|}
\hline
\textbf{Method} & \textbf{Year} & \textbf{Model Type} & \textbf{EER} & \textbf{Enrollment} & \textbf{Cross-User?} \\
\hline
Manhattan (scaled)~\cite{b4} & 2009 & Statistical & 9.6\% & 200 samples & No \\
Nearest Neighbor (Mahalanobis)~\cite{b4} & 2009 & Statistical & 10.2\% & 200 samples & No \\
Random Forest~\cite{b5} & 2013 & ML Ensemble & 8.5\% & 100 samples & Limited \\
TypeNet~\cite{b13} & 2021 & LSTM & 4.8\% & 50 samples & Limited \\
CNN-RNN~\cite{b8} & 2020 & Deep Learning & 5.2\% & 50 samples & No \\
\hline
\textbf{Ours (Encoder only)} & 2024 & 1D-CNN + SupCon & 17.7\%* & 5--10 samples & \textbf{Yes} \\
\textbf{Ours (Encoder + Adapter)} & 2024 & 1D-CNN + Adapter & \textbf{2.6\%} & 5--10 samples & No \\
\hline
\multicolumn{6}{l}{\footnotesize *Cross-user EER on unseen users (different evaluation protocol from prior work)} \\
\end{tabular}
\end{table*}

Our encoder-only approach is the first to demonstrate genuine cross-user generalization on the CMU benchmark, enabling verification of completely unseen users with reasonable accuracy. Combined with per-user adapters, our method achieves state-of-the-art in-distribution EER of 2.6\% while requiring only 5--10 enrollment samples---5--20$\times$ fewer than prior methods.

\section{Novelty and Contributions}
\label{sec:novelty}

\subsection{First Application of Supervised Contrastive Learning to Keystroke Dynamics}

Prior work on keystroke dynamics uses:
\begin{itemize}
\item Cross-entropy loss for user classification~\cite{b7,b13}
\item Triplet loss with random sampling~\cite{b21}
\item Siamese networks with pairwise training~\cite{b22}
\end{itemize}

We are the first to apply Supervised Contrastive Loss (SupCon) to keystroke dynamics. SupCon offers several advantages over these alternatives:
\begin{itemize}
\item Uses all same-user samples as positives simultaneously, not just pairs or triplets
\item Trains on mini-batches with efficient GPU utilization
\item Produces embeddings that generalize to unseen users without retraining
\item Temperature-controlled clustering that balances tightness and optimization stability
\end{itemize}

\subsection{First Two-Stage Generalization + Personalization Framework}

Previous approaches either:
\begin{itemize}
\item Train a single model for all users (generalization only)
\item Train separate models per user (personalization only)
\end{itemize}

Our two-stage framework explicitly separates these concerns:
\begin{enumerate}
\item \textbf{Stage 1:} Global encoder for generalizable representations
\item \textbf{Stage 2:} Per-user adapters for personalization
\end{enumerate}

This separation enables fast enrollment with minimal samples while maintaining strong verification performance.

\subsection{Analysis of Adapter Behavior}

We provide the first systematic analysis of why per-user adapters:
\begin{itemize}
\item Improve performance for enrolled users (personalization)
\item Do not generalize to unseen users (by design)
\end{itemize}

This clarifies the appropriate evaluation protocol for adaptive biometric systems.

\section{Limitations and Future Work}
\label{sec:limitations}

\subsection{Current Limitations}

While our framework achieves strong performance, several limitations should be acknowledged:

\begin{enumerate}
\item \textbf{Fixed-Text Constraint:} Our evaluation uses fixed-text keystroke data where users type a single, known password. This controlled setting maximizes timing consistency but limits real-world applicability. Free-text keystroke dynamics, where users type arbitrary content, requires handling variable-length sequences and is significantly more challenging due to vocabulary and context variations.

\item \textbf{Single Dataset:} We evaluate exclusively on the CMU Benchmark with 51 users. While this is the most widely-used keystroke dataset, it was collected under controlled laboratory conditions that may not reflect real-world variability. Larger and more diverse datasets (e.g., Aalto Mobile Keystroke~\cite{b32} with 37,000 users) are needed to validate scalability.

\item \textbf{Session Drift and Temporal Stability:} The CMU dataset was collected over a relatively short period (weeks). Long-term changes in typing behavior due to fatigue, injury, aging, or device changes are not captured. Practical deployments may require periodic re-enrollment or continuous adapter updates to maintain accuracy.

\item \textbf{Adversarial Robustness:} We do not evaluate robustness to adversarial attacks such as:
\begin{itemize}
\item Keystroke injection attacks that forge timing patterns
\item Mimicry attacks where adversaries practice imitating a target's typing rhythm
\item Model inversion attacks that attempt to reconstruct typing patterns from templates
\end{itemize}

\item \textbf{Hardware Dependence:} Keystroke timings can vary across different keyboards due to key travel distance, mechanical vs. membrane switches, and scan rate. Our evaluation uses a single keyboard type; cross-device generalization requires further investigation.

\item \textbf{Population Bias:} The CMU dataset has limited demographic diversity (primarily university students). Performance may vary across different age groups, typing proficiencies, and cultural backgrounds.
\end{enumerate}

\subsection{Future Research Directions}

\begin{enumerate}
\item \textbf{Free-Text Extension:} Extend the encoder to handle variable-length keystroke sequences using attention mechanisms or Transformer architectures. N-gram based representations could provide fixed-size inputs for arbitrary text while preserving discriminative information.

\item \textbf{Transformer Architecture:} Replace the 1D-CNN with Transformer~\cite{b33} to model long-range key dependencies and capture relationships between non-adjacent keystrokes. Self-attention may be particularly effective for passwords with repeated characters or patterns.

\item \textbf{Federated Learning:} Train per-user adapters on-device without transmitting raw keystroke data to a central server. This preserves privacy while enabling collaborative improvements to the shared encoder through federated averaging.

\item \textbf{Active Threshold Adaptation:} Implement dynamic threshold adjustment based on recent verification history. Users with consistently high similarity scores could have tighter thresholds, while those with variable behavior could have relaxed thresholds to balance security and usability.

\item \textbf{Multi-Modal Fusion:} Combine keystroke dynamics with complementary behavioral biometrics:
\begin{itemize}
\item Mouse dynamics (movement patterns, click timing)
\item Touchscreen gestures (swipe patterns, pressure)
\item Gaze patterns during typing
\item Device sensor data (accelerometer, gyroscope)
\end{itemize}
Multi-modal fusion could improve robustness and enable cross-device authentication.

\item \textbf{Continuous Adapter Learning:} Implement online adapter updates that incrementally learn from verified sessions. This would address session drift while maintaining security through verification-gated updates.

\item \textbf{Low-Resource Deployment:} Investigate model compression techniques (quantization, pruning, knowledge distillation) for deployment on resource-constrained edge devices such as IoT systems and embedded controllers.

\item \textbf{Explainable Authentication:} Develop interpretability methods that explain why a verification attempt failed, helping legitimate users understand and correct their typing behavior.
\end{enumerate}

\section{Conclusion}
\label{sec:conclusion}

We presented a two-stage framework for keystroke dynamics authentication that combines supervised contrastive pretraining with per-user adapter fine-tuning. Our comprehensive evaluation on the CMU Keystroke Dynamics Benchmark yields the following key findings:

\begin{enumerate}
\item \textbf{Supervised Contrastive Learning Enables Cross-User Generalization:} A 1D-CNN encoder trained with SupCon achieves AUC = 0.90 and EER = 17.7\% for cross-user verification on completely unseen users, substantially outperforming Random Forest (AUC = 0.80, EER = 26.3\%) and XGBoost (AUC = 0.66, EER = 39.4\%). This is the first demonstration of genuine cross-user generalization in keystroke dynamics.

\item \textbf{Lightweight Adapters Enable Efficient Personalization:} Per-user adapters with only 16,576 parameters, trained on 5--10 enrollment samples in under 2 seconds, reduce EER from 3.3\% to 2.6\%---a 21\% relative improvement. The adapter architecture achieves state-of-the-art in-distribution performance while requiring 5--20$\times$ fewer enrollment samples than prior methods.

\item \textbf{Clear Separation of Generalization and Personalization:} Our analysis reveals that encoders and adapters serve fundamentally different roles. The encoder provides generalizable representations that transfer to unseen users; adapters provide user-specific personalization that does not transfer. This understanding is crucial for proper evaluation of adaptive biometric systems.

\item \textbf{Practical Deployment Characteristics:} The framework supports real-time verification ($<$4 ms on CPU), modest storage requirements ($\sim$67 KB per user), and seamless transition from generalization mode (pre-enrollment) to personalization mode (post-enrollment).
\end{enumerate}

Our key insight is that per-user adapters are personalization modules effective only after enrollment, not generalization mechanisms. Conflating these roles leads to inappropriate evaluation protocols and misleading conclusions. We hope this work provides a foundation for future research in adaptive behavioral biometrics.

\section*{Data Availability}
The CMU Keystroke Dynamics Benchmark dataset used in this study is publicly available at \url{https://www.cs.cmu.edu/~keystroke/}. Our code, pretrained encoder, and per-user adapters are available at [repository URL] for reproducibility.

\section*{Acknowledgment}
The authors thank the Carnegie Mellon University for creating and maintaining the Keystroke Dynamics Benchmark dataset. We also acknowledge the contributions of the deep metric learning and adapter communities whose foundational work enabled this research.

\begin{thebibliography}{00}

\bibitem{b1} F. Monrose and A. D. Rubin, ``Keystroke dynamics as a biometric for authentication,'' \emph{Future Generation Computer Systems}, vol. 16, no. 4, pp. 351--359, 2000.

\bibitem{b2} R. Gaines, W. Lisowski, S. Press, and N. Shapiro, ``Authentication by keystroke timing: Some preliminary results,'' RAND Corporation, Tech. Rep. R-2526-NSF, 1980.

\bibitem{b3} A. Alsultan and K. Warwick, ``A review of keystroke dynamics biometrics approaches,'' \emph{Int. J. Computer and Information Technology}, vol. 3, no. 3, pp. 544--549, 2013.

\bibitem{b4} K. S. Killourhy and R. A. Maxion, ``Comparing anomaly-detection algorithms for keystroke dynamics,'' in \emph{Proc. IEEE/IFIP Int. Conf. Dependable Systems and Networks}, 2009, pp. 125--134.

\bibitem{b5} P. S. Teh, A. B. J. Teoh, and S. Yue, ``A survey of keystroke dynamics biometrics,'' \emph{The Scientific World Journal}, vol. 2013, pp. 1--24, 2013.

\bibitem{b6} S. P. Banerjee and D. L. Woodard, ``Biometric authentication and identification using keystroke dynamics: A survey,'' \emph{J. Pattern Recognition Research}, vol. 7, no. 1, pp. 116--139, 2012.

\bibitem{b7} Y. Deng and Y. Zhong, ``Keystroke dynamics advances using deep neural networks,'' in \emph{Recent Advances in Keystroke Dynamics Biometrics}, 2013, pp. 59--70.

\bibitem{b8} X. Lu, S. Zhang, P. Hui, and P. Lio, ``Continuous authentication by free-text keystroke based on CNN and RNN,'' \emph{Computers and Security}, vol. 96, pp. 101861, 2020.

\bibitem{b9} A. Morales, J. Fierrez, R. Tolosana, et al., ``Keystroke biometrics ongoing competition,'' \emph{IEEE Access}, vol. 4, pp. 7736--7746, 2016.

\bibitem{b10} P. Khosla, P. Teterwak, C. Wang, et al., ``Supervised contrastive learning,'' in \emph{Advances in Neural Information Processing Systems}, vol. 33, 2020, pp. 18661--18673.

\bibitem{b11} F. Bergadano, D. Gunetti, and C. Picardi, ``User authentication through keystroke dynamics,'' \emph{ACM Trans. Information and System Security}, vol. 5, no. 4, pp. 367--397, 2002.

\bibitem{b12} S. Cho, C. Han, D. H. Han, and H. I. Kim, ``Web-based keystroke dynamics identity verification using neural network,'' \emph{J. Organizational Computing and Electronic Commerce}, vol. 10, no. 4, pp. 295--307, 2000.

\bibitem{b13} A. Acien, A. Morales, J. Fierrez, R. Vera-Rodriguez, and J. Hernandez-Ortega, ``TypeNet: Deep learning keystroke biometrics,'' \emph{IEEE Trans. Biometrics, Behavior, and Identity Science}, vol. 4, no. 1, pp. 57--70, 2021.

\bibitem{b14} S. Roy, U. Roy, and D. D. Sinha, ``A deep learning approach for keystroke dynamics,'' \emph{Expert Systems with Applications}, vol. 191, pp. 116222, 2022.

\bibitem{b15} T. Chen, S. Kornblith, M. Norouzi, and G. Hinton, ``A simple framework for contrastive learning of visual representations,'' in \emph{Proc. ICML}, 2020, pp. 1597--1607.

\bibitem{b16} F. Schroff, D. Kalenichenko, and J. Philbin, ``FaceNet: A unified embedding for face recognition and clustering,'' in \emph{Proc. IEEE CVPR}, 2015, pp. 815--823.

\bibitem{b17} A. K. Jain, K. Nandakumar, and A. Ross, ``50 years of biometric research: Accomplishments, challenges, and opportunities,'' \emph{Pattern Recognition Letters}, vol. 79, pp. 80--105, 2016.

\bibitem{b18} J. Snell, K. Swersky, and R. Zemel, ``Prototypical networks for few-shot learning,'' in \emph{Advances in Neural Information Processing Systems}, vol. 30, 2017.

\bibitem{b19} G. Koch, R. Zemel, and R. Salakhutdinov, ``Siamese neural networks for one-shot image recognition,'' in \emph{ICML Deep Learning Workshop}, 2015.

\bibitem{b20} N. Houlsby, A. Giurgiu, S. Jastrzebski, et al., ``Parameter-efficient transfer learning for NLP,'' in \emph{Proc. ICML}, 2019, pp. 2790--2799.

\bibitem{b21} P. H. Pisani and A. C. Lorena, ``Adaptive approaches for keystroke dynamics,'' \emph{Expert Systems with Applications}, vol. 126, pp. 153--167, 2019.

\bibitem{b22} R. Giot, M. El-Abed, and C. Rosenberger, ``Keystroke dynamics authentication,'' in \emph{Proc. Int. Symp. Collaborative Technologies and Systems}, 2009, pp. 172--179.

\bibitem{b23} Y. Movshovitz-Attias, A. Toshev, T. K. Leung, S. Ioffe, and S. Singh, ``No fuss distance metric learning using proxies,'' in \emph{Proc. IEEE ICCV}, 2017, pp. 360--368.

\bibitem{b24} J. Deng, J. Guo, N. Xue, and S. Zafeiriou, ``ArcFace: Additive angular margin loss for deep face recognition,'' in \emph{Proc. IEEE CVPR}, 2019, pp. 4690--4699.

\bibitem{b25} Y. Wen, K. Zhang, Z. Li, and Y. Qiao, ``A discriminative feature learning approach for deep face recognition,'' in \emph{Proc. ECCV}, 2016, pp. 499--515.

\bibitem{b26} X. Wang, X. Han, W. Huang, D. Dong, and M. R. Scott, ``Multi-similarity loss with general pair weighting for deep metric learning,'' in \emph{Proc. IEEE CVPR}, 2019, pp. 5022--5030.

\bibitem{b27} O. Vinyals, C. Blundell, T. Lillicrap, K. Kavukcuoglu, and D. Wierstra, ``Matching networks for one shot learning,'' in \emph{Advances in Neural Information Processing Systems}, vol. 29, 2016.

\bibitem{b28} J. Chen, S. H. Kao, H. He, et al., ``Run, don't walk: Chasing higher FLOPS for faster neural networks,'' in \emph{Proc. IEEE CVPR}, 2023, pp. 12021--12031.

\bibitem{b29} E. J. Hu, Y. Shen, P. Wallis, et al., ``LoRA: Low-rank adaptation of large language models,'' in \emph{Proc. ICLR}, 2022.

\bibitem{b30} N. K. Ratha, J. H. Connell, and R. M. Bolle, ``Enhancing security and privacy in biometrics-based authentication systems,'' \emph{IBM Systems Journal}, vol. 40, no. 3, pp. 614--634, 2001.

\bibitem{b31} A. Juels and M. Wattenberg, ``A fuzzy commitment scheme,'' in \emph{Proc. ACM CCS}, 1999, pp. 28--36.

\bibitem{b32} A. Acien, A. Morales, J. V. Monaco, R. Vera-Rodriguez, and J. Fierrez, ``TypeNet: Scaling up keystroke biometrics,'' in \emph{Proc. IEEE IJCB}, 2020, pp. 1--7.

\bibitem{b33} A. Vaswani, N. Shazeer, N. Parmar, et al., ``Attention is all you need,'' in \emph{Advances in Neural Information Processing Systems}, vol. 30, 2017, pp. 5998--6008.

\end{thebibliography}

\begin{IEEEbiography}[{\includegraphics[width=1in,height=1.25in,clip,keepaspectratio]{white bg me.jpg}}]{\textbf{Atul Kumar Mishra}}
is an undergraduate student in the Computer Science and Engineering (Cyber Security) program at R.V. College of Engineering, Bengaluru, India. His academic interests include computer vision, embedded systems, and cybersecurity. His technical skill set includes Python, C++, Java, OpenCV, PyTorch, Flask, and MQTT, with a strong focus on applying these tools to practical and research-oriented challenges.
\end{IEEEbiography}

\begin{IEEEbiography}[{\includegraphics[width=1in,height=1.25in,clip,keepaspectratio]{harsh.jpeg}}]{\textbf{Harsh Sagar}}
is an undergraduate student in the Department of Computer Science and Engineering at R.V. College of Engineering, Bengaluru, India. His interests include distributed systems, backend development, and applied cybersecurity. He is particularly interested in building scalable systems, exploring system performance optimization, and applying secure software engineering practices to real-world applications.
\end{IEEEbiography}

\begin{IEEEbiography}[{\includegraphics[width=1in,height=1.25in,clip,keepaspectratio]{amit.jpeg}}]{\textbf{Amitansh Gupta}}
is a cybersecurity enthusiast with a practical skill set in building and securing software systems. He works with the MERN stack and has experience with Python, C++, C, Flask, PHP, and MongoDB. His interests include machine learning, convolutional neural networks, and modern encryption methods, which he applies to understanding and protecting digital systems. He enjoys learning new tools and solving technical problems.
\end{IEEEbiography}

\begin{IEEEbiography}[{\includegraphics[width=1in,height=1.25in,clip,keepaspectratio]{suma.png}}]{\textbf{Dr.\ Suma B}}
is currently an academic and researcher with expertise in Artificial Intelligence and Machine Learning. Her primary research focus lies in artificial intelligence, machine learning, and data mining, with allied interests in data hiding techniques and algorithm analysis. Her work emphasizes developing efficient algorithms and intelligent systems for data-driven applications.
\end{IEEEbiography}




\end{document}
